\chapter{Background}
\label{chap:background}

This chapter positions the project against existing work and
establishes the concepts the rest of the report depends on.
Section~\ref{sec:background-problem} defines the magazine translation
problem and argues that layout preservation is part of translation
correctness, with particular attention to the typographic divergence
between Latin and Chinese. Section~\ref{sec:background-routes} surveys
the three automatic routes to layout-preserving translation and
justifies this project's choice of the PDF-native route.
Section~\ref{sec:background-ir} examines intermediate representations
for documents through three case studies, and
Section~\ref{sec:background-babeldoc} analyses BabelDoc, the system
this project extends, in depth --- including the boundary between its
work and this project's. Section~\ref{sec:background-agentic}
introduces the ReAct pattern and defines the parameter-bounded
autonomy contract adopted here;
Section~\ref{sec:background-hitl} reviews constrained and
human-in-the-loop translation; and
Section~\ref{sec:background-eval} reviews evaluation instruments for
translated documents, from sentence-level metrics to document-level
coherence and geometric layout quality.

\section{The Magazine Translation Problem and Why Layout Matters}
\label{sec:background-problem}

\subsection{Layout as Part of Translation Correctness}

A magazine page is not a container for text; it is an instruction for
how to read. Column order determines sentence order across column
breaks. Visual hierarchy --- headline, standfirst, body, caption ---
tells the reader what is content and what is apparatus. Devices such
as pull quotes duplicate body text for emphasis, drop caps mark article
openings, and sidebars carry content that is related to, but must not
be interleaved with, the main article. When several articles share a
page, or one article spans several pages, only the layout distinguishes
one text from another and indicates where each continues.

Layout preservation is therefore part of translation
\emph{correctness}, not an aesthetic afterthought. If a translated
page breaks column order, the reader receives sentences in the wrong
sequence; if it collapses hierarchy, apparatus and content become
indistinguishable; if it merges blocks across article boundaries, the
output interleaves two texts that were never one. Each of these is a
meaning-level failure produced entirely at the layout level, and none
of them is detectable in the translated text alone. This is the
central reason the present project treats the rendered page, not the
extracted text, as the unit of evaluation
(Section~\ref{sec:background-eval}).

\subsection{Typographic Divergence between Latin and Chinese}
\label{sec:background-typography}

English-to-Chinese translation makes the layout problem harder in ways
that are specific to the script pair, and that reward being stated
precisely \cite{w3c_clreq}.

First, the two scripts justify text by different mechanisms. Latin
typography sets words separated by spaces, and justifies a paragraph
by distributing adjustment across inter-word spaces, letter spacing,
and hyphenation. Chinese typography sets a continuous run of
predominantly full-width characters with no word delimiters:
justification cannot lean on inter-word spaces that do not exist;
the primary fine adjustment is the compression of full-width
punctuation, with inter-character spacing admitted only as a
secondary method \cite{w3c_clreq}. The \emph{inter-word} axis on
which Latin justification chiefly relies is thus absent entirely,
and the remaining axes are conventionally used far more sparingly.
%% VERIFIED V01 2026-08-14 https://www.w3.org/TR/clreq/#line_adjustment
%% clreq §6.2.2: line adjustment applies punctuation-width adjustment
%% first and adds space between Han characters only when that is
%% insufficient.
The main axis of variation that Latin compositors use to fit text into
a fixed geometry is thus absent on the Chinese side; where an English
paragraph can be coaxed into a tight box by re-spacing, the Chinese
translation of that paragraph largely either fits or does not.

Second, line breaking follows different rules. Chinese permits a break
between almost any two characters, subject to prohibition rules that
forbid certain punctuation at line starts and ends \cite{w3c_clreq},
and it does not hyphenate; English breaks only at spaces and
hyphenation points. Line-break opportunity is therefore denser in
Chinese but constrained differently, and naive re-use of the source's
line geometry produces violations of the target script's conventions.

Third, the visual footprint of the same content changes. A Chinese
rendering typically uses far fewer characters than its English source,
but each character occupies a full em square; the translated block is
therefore generally not the same size or shape as the space vacated by
the source text.
%% VERIFIED V02 2026-08-14 resolved qualitative; optional corpus
%% footprint ratio deferred to evaluation chapter (owner note)
Reinserting translated text into the exact source boxes consequently
tends to produce either large blank areas or overflow --- the failure
mode observed in the pilot studies of Section~\ref{sec:background-routes} --- and
this, together with the justification asymmetry above, is the concrete
motivation for treating text placement as a reflow problem over
article-level regions rather than a box-filling problem over
block-level regions (Chapter~\ref{chap:method}).

Throughout this report, ``Chinese'' refers to Simplified Chinese, and
English-to-Chinese magazine translation is the primary studied
direction; the corpus and direction choices are justified in
Section~\ref{sec:corpus}.
%% NOTE: this sentence absorbs the scoping material that the interim
%% report carried in Section 1.1 (which the second marker judged too
%% early there); keep Chapter 1 down to one forward reference.
% =====================================================================
% Section 2.2 — Automatic Routes to Layout-Preserving Translation
% (draft v1)
% Three-route taxonomy replacing interim sections 2.1 (partly) and 2.2.
% Conventions as before. Figures reuse the interim filenames.
% MMT paragraph (mmt_survey_2024, futeral2023commute) deliberately
% omitted pending the architecture decision on visual context in
% translation; if the final translator consumes page images, add one
% sentence at the end of 2.2.2 rather than restoring the old paragraph.
% =====================================================================

\section{Automatic Routes to Layout-Preserving Translation}
\label{sec:background-routes}

Automatic approaches that preserve layout while translating fall into
three routes, distinguished by what they take as input and what they
make explicit in between: commercial end-to-end systems, academic work
on document \emph{image} translation, and PDF-native pipelines that
consume the document's own text layer and geometry. This section
surveys all three, explains why this project builds on the third, and
identifies the limitation the three routes share.

\subsection{Commercial Document Translation Systems}

DeepL Document Translation translates uploaded documents while aiming
to preserve their layout and design \cite{deepl_document_translation}.
Its engineering discussion of PDF translation is unusually informative
about mechanism: the pipeline reconstructs the document after
translation, and the team describes using page position and
bounding-box overlap as signals of layout quality
\cite{deepl_document_workflow} --- an early industrial precedent for
the geometric evaluation measures discussed in
Section~\ref{sec:background-layoutmetrics}. In the pilot comparison
run for this project (Figure~\ref{fig:deepl-pilot}), DeepL produced
relatively fluent Chinese and preserved the broad structure of the
page, but when translated text did not fit the source geometry it
moved the excess onto an additional, sparsely filled page --- an
acceptable compromise for office documents, but not for a magazine,
where the page grid is part of the product.

\begin{figure}[tbp]
    \centering
    \includegraphics[width=\linewidth]{\detokenize{figures/deepl pilot.png}}
    \caption{DeepL Document Translation preserves part of the page
    structure, but may move overflow content to an additional sparse
    page.}
    \label{fig:deepl-pilot}
\end{figure}

Google Cloud Document Translation likewise aims to preserve original
formatting and layout \cite{google_document_translation}. In the pilot
(Figure~\ref{fig:google-pilot}) it was fast and retained image
placement and broad page geometry, but text segmentation followed
visual lines rather than semantic paragraphs: multi-column body text
and pull quotes were fragmented into line-level units, which damages
both readability and --- because each fragment is translated with
truncated context --- translation accuracy itself.

\begin{figure}[!htbp]
    \centering
    \includegraphics[width=\linewidth]{\detokenize{figures/google pilot.png}}
    \caption{Google Document Translation can split paragraphs into
    line-level fragments, degrading both layout and translation
    accuracy.}
    \label{fig:google-pilot}
\end{figure}

A third commercial-adjacent route generates the translated page as an
image. GPT-based visual translation (Figure~\ref{fig:gpt-pilot})
produced natural-looking, magazine-styled Chinese pages, but the
output behaves as a picture rather than a document: characters may be
blurred or malformed --- consistent with prior findings that image
generation models struggle to render well-formed glyphs
\cite{liu2022characteraware} --- and the text is not reliably
editable, searchable, or traceable to its source.

\begin{figure}[!htbp]
    \centering
    \includegraphics[width=0.70\linewidth]{\detokenize{figures/gpt pilot.png}}
    \caption{GPT-based visual translation can look plausible, but
    distorted or blurred characters are not standard Chinese glyphs,
    and the output is an image rather than structured text.}
    \label{fig:gpt-pilot}
\end{figure}

These pilots are motivating observations, not a benchmark: each system
was run in its default configuration on a small number of pages, and
the observations identify failure modes rather than measure rates.
%% NOTE: this framing sentence survived from the interim report and
%% was not criticised; keep it.

\subsection{The Image Route: Document Image Translation}
\label{sec:background-routes-image}

Academic work has concentrated on the harder input condition where the
document is only available as an image. The classical pipeline runs
OCR and then machine translation, but OCR output fragments paragraphs
into lines, misses decorative or image-internal text, and propagates
its recognition errors into every downstream stage --- a weakness
identified explicitly by Kim et al., whose Donut model removes OCR
entirely and maps the page image to structured output end-to-end,
improving over OCR-dependent pipelines on information-extraction
benchmarks (on CORD, 84.1 vs.\ 78.9 F1) while avoiding per-language
OCR engineering
\cite{kim2021donut}.
%% VERIFIED V03 2026-08-14 https://arxiv.org/abs/2111.15664 (Tab. 2;
%% claim scoped to the information-extraction benchmarks per owner
%% adjudication)
In the same spirit, multimodal document models such as LayoutLMv3
jointly encode text and image with unified masking objectives and
word--patch alignment, improving document understanding tasks over
text-only encoding: its base model reaches 90.29 F1 on FUNSD and
96.56 F1 on CORD, against 66.48 and 93.54 for the comparably sized
text-only RoBERTa \cite{huang2022layoutlmv3}.
%% VERIFIED V04 2026-08-13 https://arxiv.org/pdf/2204.08387 (Table 1)

Document image translation (DIT) work builds translation directly into
this setting. Translatotron-V translates text inside images end-to-end
and introduces Structure-BLEU, which couples translation scoring to
spatial correspondence
(Section~\ref{sec:background-layoutmetrics})
\cite{lan2024translatotron}, matching the strongest cascaded
OCR--translate--render pipeline on German--English in-image
translation (15.39 against 15.37 BLEU) with roughly 71\% of its
parameters. QRDIT addresses whole-page,
complex-layout documents with a query--response framework that
concentrates translation on relevant regions \cite{zhang2025query},
improving the previous state of the art (LayoutDIT) by 2.14 average
BLEU across the three English--Chinese DITrans domains (26.84 against
24.70), and ZoomDIT recovers document-level coherence from fragmented OCR
words through a fine-to-coarse zoom-out network \cite{zhang2025chaotic},
improving the same baseline by 2.27 average BLEU on English--Chinese
DIT700K (38.72 against 36.45).
%% VERIFIED V05 2026-08-13 https://aclanthology.org/2024.findings-acl.325/ ;
%% https://aclanthology.org/2025.findings-acl.372/ (Tab. 3) ;
%% https://aclanthology.org/2025.coling-main.723/ (Tab. 3)
Collectively this line establishes techniques for exactly the page
phenomena magazines exhibit --- dense mixed layouts, region-dependent
translation, fragmented text recovery --- but always under the
assumption that the source text must be recovered from pixels.
% NOTE: if the final architecture consumes page images during
% translation, add ONE sentence here on multimodal context (restore
% keys mmt_survey_2024 / futeral2023commute); do not restore the old
% MMT paragraph.

\subsection{The PDF-Native Route}
\label{sec:background-routes-pdfnative}

When the document is born digital, its exact text, fonts, and geometry
are already present in the PDF, and a third route consumes them
directly. BabelDoc \cite{qi2026babeldoc} is the representative and,
to the knowledge gathered for this report, the only open-source system
on this route with a published system description: it parses the PDF
into an explicit intermediate representation, translates over that
representation with document-level operations, and re-anchors the
translated content into the source geometry. Because this project is
built as an extension of BabelDoc, the system is examined in depth in
Section~\ref{sec:background-babeldoc}, and its representation is compared with
other document IRs in Section~\ref{sec:background-ir}.

\subsection{Route Choice, Scope, and the Shared Gap}
\label{sec:background-routes-gap}

\paragraph{Route choice.} Contemporary magazines are produced with
digital layout tools, and their PDFs carry a complete text layer. For
such documents the image route answers a harder question than the one
being asked: recovering text from pixels that the file already
contains verbatim, at the cost of OCR error propagation
\cite{kim2021donut} and the loss of exact fonts and geometry. This
project therefore takes the PDF-native route. An early prototype of
this project rasterised pages and applied OCR before this decision was
revisited; the change and its rationale are recorded in
Chapter~\ref{chap:method}.
%% NOTE: one honest sentence about the early OCR prototype
%% pre-empts the question the second marker asked about rasterisation,
%% without re-arguing it here.

\paragraph{Scope.} The PDF-native route defines this project's
applicability boundary: a PDF is not guaranteed to contain an
extractable text layer, and scanned or legacy magazine issues do not.
For those inputs the image route of
Section~\ref{sec:background-routes-image} is the appropriate technology, and
extending the present system with an image-route front end is future
work. All corpora used in this project are born-digital
(Section~\ref{sec:corpus}), and the generalisation claims of
Chapter~\ref{chap:eval} are made within this boundary.

\paragraph{The shared gap.} Across all three routes, processing is
scoped to the page or below: blocks are detected per page, translated
per page, and re-placed per page. None of the routes carries an
explicit unit for the \emph{article} --- the text that readers
actually follow, which routinely spans pages and shares pages with
other articles. The consequences are concrete: a sentence split by a
page break is translated as two fragments that cannot see each other,
and terminology drifts between pages of the same article. Direct
computational work on magazine translation specifically remains
scarce, and the article-level gap identified here is, to the best of
the review conducted for this project, unaddressed on any of the three
routes; closing it is the central technical contribution of this
project (Chapter~\ref{chap:method}), and the failure it causes is
demonstrated empirically in Section~\ref{sec:eval-midbreak}.
% =====================================================================
% Section 2.3 — Intermediate Representations for Documents (draft v1)
% Case-study section requested explicitly by the second marker.
% Same conventions as section_2_4_babeldoc.tex:
%  - %% VERIFY marks facts to re-check against primary sources.
%  - \ref placeholders: \ref{sec:background-routes} (2.2), \ref{sec:background-babeldoc} (2.4),
%    \ref{chap:method}, \ref{sec:method-two-layer}.
%  - Bib keys reused from interim report: docling2025, smoldocling2025;
%    new key qi2026babeldoc (defined with Section 2.4 draft).
% =====================================================================

\section{Intermediate Representations for Documents}
\label{sec:background-ir}

The routes surveyed in Section~\ref{sec:background-routes} differ in how much of
the document they make explicit between input and output. This section
examines the design space of explicit document representations, because
the representation chosen determines what a translation pipeline can
know, check, and revise. The term \emph{intermediate representation}
(IR) is used here in the same sense as in compiler construction: a
structured, inspectable encoding of the artefact being transformed,
which decouples analysis from generation and allows each stage to be
validated independently. A pipeline that passes plain text between
stages cannot answer questions such as ``which source block does this
translated sentence come from?'' or ``was every extracted block
rendered?''; a pipeline built around an explicit IR can. Three recent
systems illustrate three distinct points in this design space, and each
is examined below for what it represents, what it demonstrably does
well, and what it lacks for the magazine translation task.

\subsection{Case Study: DoclingDocument}

Docling is an open-source document conversion toolkit built around a
unified representation called \texttt{DoclingDocument}, defined as a
typed (Pydantic) data model \cite{docling2025}. The representation
expresses text, tables, pictures, captions and lists; document
hierarchy through sections and groups; an explicit distinction between
body content and \emph{furniture} such as headers and footers; layout
information as bounding boxes for all items; and provenance information
such as page numbers and document origin \cite{docling2025}.
Construction is performed by specialised models --- a layout detector
derived from RT-DETR and trained on DocLayNet, and the TableFormer
table-structure model --- and the resulting document can be traversed
in reading order, serialised losslessly to JSON, or exported to
Markdown and HTML \cite{docling2025}.
%% VERIFIED V06 2026-08-13 https://arxiv.org/abs/2501.17887 (toolkit paper
%% states the layout model is "derived from RT-DETR ... re-trained on
%% DocLayNet" and names TableFormer; DoclingDocument is defined there as a
%% Pydantic model with furniture and provenance)

Two design decisions in \texttt{DoclingDocument} are directly relevant
to this project. The furniture distinction shows the value of
separating repeating page apparatus from article content --- the same
separation magazine pages require for rails, folios, and running heads.
The provenance model shows that a representation can carry ``where did
this come from'' alongside ``what is this'', which is the precondition
for any auditable pipeline. What the representation lacks, reasonably
given its purpose, is anything specific to transformation: it records
what a document \emph{is}, not what is being \emph{done to it}. There
is no translation state, no source-to-target linkage, and its hierarchy
assumes an essentially linear document of sections --- it has no unit
for several independent articles sharing one page, nor for one article
spanning several pages.

\subsection{Case Study: DocTags}

SmolDocling takes a different position: instead of assembling a
representation from specialised detectors, a single ultra-compact
vision--language model (256M parameters, combining a SigLIP visual
backbone with a SmolLM-2 language backbone) reads a page image and
autoregressively generates the full page representation in a markup
format called DocTags \cite{smoldocling2025}. DocTags defines a
structured vocabulary that separates textual content from document
structure, covering captions, charts, forms, code, equations, tables,
footnotes, lists, page headers and footers, section headings, and body
text --- fourteen layout classes in total --- with each element
annotated with its bounding box, and OCR performed within elements as
part of the same generation \cite{smoldocling2025}. The result is
competitive with substantially larger models: on full-page text
recognition over DocLayNet the 256M-parameter model reports an F1 of
0.80 and an edit distance of 0.48, against 0.72 and 0.56 for the
7B-parameter Qwen2.5-VL \cite{smoldocling2025}.
%% VERIFIED V07 2026-08-13 https://arxiv.org/abs/2503.11576 (Table 2;
%% the "27 times larger" figure in the paper is a parameter-count
%% comparison and is deliberately not used here as a quality margin)

For this project, DocTags demonstrates two things. First, that a typed
role vocabulary --- a closed set of named element classes --- makes a
page representation both compact and checkable; this idea recurs in the
page-kind vocabulary of Chapter~\ref{chap:method}. Second, and by way
of contrast, it illustrates the cost of an end-to-end generative
representation for a translation workflow: DocTags is produced one page
at a time from a page image, so the representation is inherently
page-scoped, and its content passes through autoregressive generation
rather than extraction from the source. For a born-digital magazine
PDF, where the exact source text is available, regenerating that text
through a model reintroduces precisely the fidelity risk that
PDF-native extraction avoids (Section~\ref{sec:background-routes}).

\subsection{Case Study: The BabelDoc IL}

BabelDoc's Intermediate Language, described in detail in
Section~\ref{sec:background-babeldoc}, occupies a third position: unlike the two
representations above, it is designed \emph{for} transformation rather
than for conversion. The IL records character-level text, style, and
detected layout regions from the source PDF, and the pipeline's
document-level operations --- terminology extraction,
glossary-constrained generation, formula placeholdering --- operate
over this state before an adaptive typesetting engine re-anchors the
translated content into the source geometry \cite{qi2026babeldoc}. The
IL therefore does carry transformation state, and it is the only
representation of the three in which ``source block'' and ``translated
block'' are linked objects. Its structural ceiling, established in
Section~\ref{sec:background-babeldoc}, is that this state stops at the block and
page level: there is no unit between block and document, so article
identity, cross-page continuity, and article-scoped context cannot be
represented.

\subsection{Synthesis}

Table~\ref{tab:ir-comparison} summarises the three representations
against the requirements of layout-preserving magazine translation.
Each contributes a proven idea; none covers the requirements alone.

\begin{table}[tb]
\centering
\small
\begin{tabular}{p{0.26\linewidth}p{0.20\linewidth}p{0.20\linewidth}p{0.22\linewidth}}
\hline
 & \textbf{DoclingDocument} & \textbf{DocTags} & \textbf{BabelDoc IL} \\
\hline
Primary purpose & conversion & end-to-end conversion & translation \\
Construction & specialised detectors & generative VLM & PDF-native parsing \\
Structural units & section hierarchy, furniture & typed page elements & blocks, pages \\
Cross-page semantics & none & none (page-scoped) & context sharing only (\S\ref{sec:background-babeldoc}) \\
Transformation state & none & none & source--target linkage \\
Human adjudication & none & none & none \\
\hline
\end{tabular}
\caption{Document representations compared against the requirements of
layout-preserving magazine translation. All three lack an article-level
unit and human-adjudicated semantics; only the BabelDoc IL carries
translation state.}
\label{tab:ir-comparison}
\end{table}
%% VERIFIED V08 2026-08-14 aligned with the Section 2.4 wording adopted
%% under E1: the IL shares context across a page boundary but does not
%% reconstruct a cross-page unit.

The synthesis this project draws is as follows. From
\texttt{DoclingDocument} it takes the principles that provenance
belongs inside the representation and that page apparatus should be
structurally separated from content. From DocTags it takes the value of
a closed, typed vocabulary of element and page roles. From the BabelDoc
IL it takes the translation-oriented core: an inspectable state in
which source and target remain linked through rendering. What none of
the three provides --- an authoritative unit for the article, spanning
pages and carrying translation-boundary semantics, together with a
mechanism for human-adjudicated metadata to take precedence over
machine-extracted structure --- is what the two-layer representation
developed in Chapter~\ref{chap:method} adds: an article layer whose
paragraph-level chain identifiers are the single authority for
translation-unit boundaries, and a page layer whose page kinds are
consumed only as policy priors, never as content authority.
%% NOTE: keep the final sentence structurally parallel to the boundary
%% statement in Section 2.4 --- the two passages are the ``existing
%% work ends / my work begins'' pair for representations and for the
%% base system respectively.
% =====================================================================
% Section 2.4 — BabelDoc as the Base System (draft v1)
% Style notes:
%  - Every claim about BabelDoc is sourced to the paper or project docs.
%  - %% VERIFY comments mark facts to check against the original before
%    submission (exact numbers, authorship, current README).
%  - \ref placeholders point at chapters/sections to be named later:
%    \ref{chap:method}, \ref{sec:eval-midbreak}, \ref{sec:corpus}.
%  - Suggested bib key: qi2026babeldoc (replaces any yang2026 draft key).
% =====================================================================

\section{BabelDoc as the Base System}
\label{sec:background-babeldoc}

The system developed in this project is built as an extension of
BabelDoc, an open-source, IR-based framework for layout-preserving PDF
translation \cite{qi2026babeldoc}. This section describes what BabelDoc
does and what evidence supports its design, explains why it was chosen
as the base system, and then identifies the specific limitations that
define where this project's work begins.

\subsection{Design and Evidence}

BabelDoc addresses a tension identified by its authors between
linguistic processing and layout preservation: computer-assisted
translation systems tend to discard structural metadata, while document
parsers extract structure but do not support faithful re-rendering
after translation \cite{qi2026babeldoc}. Its central design decision is
an explicit Intermediate Language (IL) that decouples visual layout
metadata from semantic content. The pipeline parses the PDF into the IL
(recording text, character-level position, style, and detected layout
regions), performs analysis and translation over the IL, and finally
re-anchors the translated content into the original page geometry
through an adaptive typesetting engine. Because translation operates on
an explicit document state rather than on raw extracted text, the
framework can support document-level operations such as terminology
extraction, glossary-constrained generation, formula placeholdering,
and a degree of cross-page context handling \cite{qi2026babeldoc}.
%% VERIFIED V09 2026-08-13 https://aclanthology.org/2026.acl-demo.25/ (the
%% paper names its stages "Multimodal Parsing", "NLP & Semantic Engine" and
%% "Adaptive Typesetting & Reconstruction"; frontend/midend/backend do not
%% occur in the paper, so the avoidance above is correct)

The framework is supported by empirical evidence rather than by design
argument alone. Its authors evaluate on a curated 200-page benchmark
spanning three domains --- scientific literature (80 pages), technical
documentation (60 pages), and international patents (60 pages) ---
using human evaluation and multimodal LLM-as-a-judge protocols
alongside geometric measures including bounding-box IoU
\cite{qi2026babeldoc}. On this benchmark, against the two baselines
they evaluate --- DeepL Document Translation and PDFMathTranslate ---
BabelDoc improves layout fidelity, visual aesthetics, and terminology
consistency: in the human evaluation it scores 4.59, 4.46, and 4.47 on
a 1--5 scale, against 3.44, 3.63, and 4.21 for DeepL and 3.29, 3.28,
and 3.34 for PDFMathTranslate, with a bounding-box IoU of 50.0\%
against 19.8\% and 48.7\% respectively
%% VERIFIED V10 2026-08-13 https://aclanthology.org/2026.acl-demo.25/ (Tab. 3)
\cite{qi2026babeldoc}. Two results deserve emphasis because they shape
how the system should be used. First, in the LLM-as-a-judge evaluation
BabelDoc's translation-precision score ties with DeepL, and in the
human evaluation DeepL leaves fewer untranslated blocks; the authors
conclude that BabelDoc's contribution lies not in replacing strong
translation engines but in providing a layout-aware, controllable
document-level pipeline around them \cite{qi2026babeldoc}. Second, the
authors attribute part of the residual error to extraction-side
omissions and rendering-stage coverage gaps \cite{qi2026babeldoc} ---
that is, to silent losses between pipeline stages. Both observations
directly motivate design choices in this project: the first locates the
research contribution at the pipeline level rather than the
translation-engine level, and the second motivates the explicit
conservation checks introduced in
Chapter~\ref{chap:method}.
%% NOTE: the conservation-invariant link here is a deliberate set-up;
%% the eval chapter should call back to this sentence.

\subsection{Why BabelDoc Was Chosen}

Three properties make BabelDoc the appropriate base among the routes
surveyed in Section~\ref{sec:background-routes}. First, it is PDF-native: it
consumes the text layer and geometry of born-digital PDFs directly,
avoiding the OCR error propagation discussed in
Section~\ref{sec:background-routes} \cite{kim2021donut}. Second, its IL is an
explicit, inspectable document state, which is a precondition for the
auditable, revisable workflow this project targets; a pipeline that
passes opaque text between stages could not support the two-layer
document representation developed in Chapter~\ref{chap:method}. Third,
it is open source with a published system description
\cite{qi2026babeldoc}, which makes a disciplined extension possible:
the boundary between upstream behaviour and this project's changes can
be stated, versioned, and audited precisely.

\subsection{Limitations for Magazine Translation}

BabelDoc was designed and validated for documents whose dominant
structure is a single linear reading flow --- papers, manuals, and
patents. Magazines violate this assumption in specific ways, and each
limitation below is grounded either in the paper, in the project's
public documentation, or in reproducible failures observed on the
corpus used in this project (Section~\ref{sec:corpus}).

First, the benchmark contains no magazine-like domain: none of its
three document classes exhibits the multi-column article layouts,
pull quotes, table-of-contents grids, sidebar boxes, or rail elements
that characterise magazine pages \cite{qi2026babeldoc}. The domain gap
is concrete rather than rhetorical: the project documentation lists
drop caps --- a standard magazine typographic device --- as
unsupported \cite{babeldoc_repo_docs}.
%% VERIFIED V11 2026-08-13 https://github.com/funstory-ai/BabelDOC (README,
%% "Known Issues": "Does not support drop caps."; accessed 2026-08-13.
%% Cited to the repository documentation, not to the paper.)

Second, and centrally for this project, the base system's cross-page
handling shares context but does not reconstruct units. Inspection of
the translation module shows the mechanism precisely: at a page
boundary, the last body-text paragraph of one page and the first of
the next --- selected by a whitelist heuristic --- are placed into the
same translation batch, so that each fragment's prompt can see the
other.\footnote{Functions \texttt{process\_cross\_page\_paragraph} and
\texttt{process\_cross\_column\_paragraph} in the upstream translation
module, inspected at the pinned upstream release v0.6.4, commit
\texttt{17480db9}; paragraph selection via the
\texttt{\_is\_body\_text\_paragraph} heuristic.} The fragments
nevertheless remain independent translation
units: each produces its own output, each output is re-anchored to its
own source box, and the typesetting stage remains strictly
page-scoped. A sentence physically split by a page break is therefore
still materialised as two outputs rendered into two boxes; whether the
splice reads as one sentence is delegated entirely to prompt-level
cooperation between two generation calls, with no structural
guarantee --- no merged unit, no redistribution of translated text
across the boundary, and no conservation relation tying the two
outputs together. The framework's claim of cross-page context handling
\cite{qi2026babeldoc} is thus accurate, but narrower than unit-level
continuity. On the corpus prepared for this project, real mid-sentence
page breaks under exactly this mechanism reproducibly caused
translation defects; the defects, their mechanism, and their repair
are analysed in Section~\ref{sec:eval-midbreak}. Notably, the official
human-translated Chinese edition of the same magazine also resorts to
mid-unit breaks in places where the extended system does not, which
suggests the problem is a genuine editorial difficulty rather than an
artefact of automatic processing (Section~\ref{sec:eval-midbreak}).
%% VERIFIED V12 2026-08-14 owner adjudication + source inspection
%% (il_translator_llm_only.py at 17480db9df92ddcb37349ce34b312335226e8ec9,
%% "Release v0.6.4", 2026-07-17; the pin is now declared in the header of
%% UPSTREAM_DIFF.md)

Third, BabelDoc has no article-level document model. Magazine pages
frequently interleave several articles, and a single article frequently
spans several pages; correct translation therefore requires knowing
which blocks belong to the same article and in what order they are
read. BabelDoc's IL records pages and blocks but carries no unit
between ``block'' and ``document'', so it cannot represent article
identity, cross-page continuity, or article-scoped terminology
context.

Fourth, terminology handling is automatic only. BabelDoc extracts
glossaries and applies them during generation \cite{qi2026babeldoc},
but there is no mechanism for a human editor to adjudicate terms,
override automatically extracted entries, or have such decisions
persist and take priority in a controlled way --- a requirement for
publication-grade translation workflows
(Section~\ref{sec:background-hitl}).

Fifth, processing strategy is fixed. The pipeline applies the same
treatment to every page, with no per-page policy selection, no bounded
configuration surface, and no audit trail recording why a given page
was processed in a given way. For heterogeneous magazine issues ---
where a table-of-contents grid, a photo spread, and a dense three-column
article demand different handling --- a single fixed strategy is
insufficient, and an unlogged adaptive one would be unauditable.

\subsection{Boundary of This Project's Work}

The limitations above define the boundary between upstream work and
this project. Everything described so far in this section is BabelDoc's
contribution. This project's contributions, developed in
Chapter~\ref{chap:method}, are: an article-level representation layered
above the IL, with paragraph-level chain identifiers as the single
authority for translation-unit boundaries; a page-level strategy layer
that maps page kinds to policy flags consumed as soft priors; a
human-in-the-loop adjudication protocol for terminology with explicit
priority semantics; a parameter-bounded controller that selects
strategies only through typed, range-bounded configuration; and a
conservation-based evaluation harness.

Two engineering disciplines keep this boundary auditable. Upstream
modifications are minimised, and every modification to BabelDoc itself
is recorded individually in a change log (\texttt{UPSTREAM\_DIFF.md})
kept under version control, so that the exact delta between the
published system and the extended system is inspectable at any commit.
All extensions live outside the upstream codebase, consume the IL
through its public structure, and route translation through a single
locked entry point, so that no behaviour of the extended system is
silently attributable to hidden upstream changes.
%% NOTE (do not delete lightly): this final paragraph is the direct
%% answer to the second marker's request to show ``where existing work
%% stops and your work begins''; if space is tight elsewhere, cut
%% elsewhere.

% ---------------------------------------------------------------------
% Suggested bibliography entry (Vancouver note style, matching
% references.bib conventions in the interim report):
%
% @misc{qi2026babeldoc,
%   note = {Qi Y, Ma X, Wang X, Wang H, Wang R. BabelDOC: better
%   layout-preserving PDF translation via intermediate representation.
%   In: Proceedings of the 64th Annual Meeting of the Association for
%   Computational Linguistics (Volume 3: System Demonstrations); 2026
%   Jul; San Diego, California, United States. Association for
%   Computational Linguistics; 2026. p. 253--262.
%   doi: 10.18653/v1/2026.acl-demo.25. Available from:
%   \url{https://aclanthology.org/2026.acl-demo.25/}.
%   arXiv:2605.10845}
% }
% %% VERIFIED V13 2026-08-13 https://aclanthology.org/2026.acl-demo.25.bib
% ---------------------------------------------------------------------
% =====================================================================
% Section 2.5 — Agentic Control and Parameter-Bounded Autonomy
% (draft v1)
% Conventions as before. \ref placeholders: \ref{chap:method},
% \ref{sec:method-controller}, \ref{sec:background-eval},
% \ref{sec:background-babeldoc}.
% Audit decisions implemented:
%  - AutoGen citation dropped; the single-controller demarcation
%    sentence stands without it (see %% NOTE below to reinstate if
%    preferred).
%  - ReAct receives its full definition HERE; Chapter 1 keeps only a
%    parenthetical forward reference "(see Section 2.5)".
%  - The interim sentence "The translation task is iterative and
%    failure-prone" is replaced by a claim scoped to THIS workflow.
% =====================================================================

\section{Agentic Control and Parameter-Bounded Autonomy}
\label{sec:background-agentic}

\subsection{ReAct and Tool-Calling}

The workflow built in this project --- parse, plan, translate, place,
render, inspect, and revise --- is iterative and failure-prone: any
stage can fail, and failures are often only observable downstream, when
a critic inspects the rendered result. A controller for such a
workflow must therefore be able to act, observe the consequences, and
revise its plan, rather than execute a fixed sequence once.

The pattern this project adopts for that controller is ReAct,
introduced by Yao et al.\ at ICLR 2023 \cite{yao2023react}. ReAct
prompts a language model to interleave \emph{reasoning traces} ---
free-text reflection on the current state of the task --- with
\emph{actions} that query tools or an environment, so that reasoning
can steer subsequent actions and observations from actions can revise
subsequent reasoning. The evidence for the pattern is substantial: on
knowledge-intensive tasks (HotpotQA, Fever), interleaving actions with
reasoning reduced the hallucination observed in reasoning-only
chain-of-thought prompting, and on interactive decision-making
benchmarks the authors report improvements over imitation- and
reinforcement-learning baselines of roughly 34 and 10 absolute
percentage points in success rate (ALFWorld and WebShop respectively),
from prompting alone with one or two in-context examples
\cite{yao2023react}.
%% VERIFIED V14 2026-08-13 https://arxiv.org/abs/2210.03629 (Sec. 1: "two or
%% even one-shot ReAct prompting is able to outperform imitation or
%% reinforcement learning methods trained with 10^3 ~ 10^5 task instances,
%% with an absolute improvement of 34% and 10% in success rates
%% respectively" --- ALFWorld and WebShop, as stated above)

The mechanism by which a model's decisions become executable is the
tool-calling pattern: the model selects a tool and emits structured
arguments; the application validates and executes the call and returns
the result as the model's next observation
\cite{openai_function_calling}. Two properties of this pattern matter
for what follows: the model never executes anything itself, and every
decision it takes is materialised as a structured, loggable object.
%% NOTE: if an academic citation is preferred alongside the vendor
%% documentation, Toolformer (Schick et al. 2023) is the standard
%% choice; add key schick2023toolformer and one clause.

In this project the controller is a single ReAct-style decision module
operating over the document state, with specialised components ---
parser, translator, layout planner, renderer, critics --- exposed to
it as tools. It is a single-controller architecture, not a multi-agent
system: no component other than the controller plans or delegates, and
the specialised components are deliberately kept as callable functions
with structured inputs and outputs.
%% NOTE: if a citation is wanted for the multi-agent contrast, the
%% dropped AutoGen reference (wu2023autogen) can be reinstated for
%% exactly this sentence and nothing else.

\subsection{From Free-Form Agency to Bounded Control}
\label{sec:background-pba}

The ReAct pattern specifies \emph{when} a controller decides ---
interleaved with observation --- but not \emph{what it is allowed to
decide}. At one extreme, the controller's action space includes
emitting free-form instructions, rules, or code that downstream stages
then obey. This maximises flexibility, and it is the prevailing style
in general-purpose agent frameworks; but for a research system it has
two costs that compound. First, auditability: when behaviour is shaped
by generated text, explaining an output requires interpreting prose,
and the effective configuration of the system at any moment has no
schema against which it can be checked. Second, reproducibility: two
runs of the same pipeline on the same input can diverge because the
controller phrased a rule differently, which makes ablation --- the
attribution of an observed improvement to a specific mechanism ---
unreliable. For a project whose evaluation rests on controlled
comparisons (Section~\ref{sec:background-eval}), an unconstrained
action space would undermine the evidence.

This project therefore adopts, and names, a narrower contract between
controller and pipeline: \emph{parameter-bounded autonomy}. The term
is defined by this project as follows. The controller retains full
autonomy over \emph{which} decisions to take and \emph{when} --- the
ReAct loop is intact --- but every decision must be expressed through
one of two typed channels: composing \emph{declarative constraints}
drawn from a closed, typed vocabulary, or adjusting \emph{numerical
parameters} each of which carries an explicitly declared admissible
range. The controller can never emit free-form rules, prompts, or
code into the pipeline. Three properties follow directly. Every
controller decision is machine-checkable: a constraint either belongs
to the vocabulary or is rejected, and a parameter either lies in its
declared range or is rejected. The system's complete decision surface
is enumerable in advance: it is the vocabulary and the ranges, both of
which are versioned configuration rather than emergent behaviour. And
runs are reproducible up to the controller's choices themselves, which
are logged as structured values and can be replayed exactly.
%% NOTE: this definition must match the method chapter's mechanics
%% word-for-word in the key terms (typed vocabulary; declared
%% admissible range; structured logging/replay). Draft the method
%% chapter's definition by copying, not paraphrasing, these clauses.

The division of labour that results is worth stating, because it
answers a natural objection. The specialised tools in this system are
unapologetically layout-specific --- drop-cap handling, column reflow,
and rail detection exist because magazines need them. Generality is
not sought in the tools but in the controller and its contract: the
controller reasons only over typed page evidence and bounded
parameters, never over publication-specific rules, so adapting the
system to a new magazine --- or, prospectively, to other layout-heavy
document families such as textbooks or illustrated manuals --- means
extending the tool set and the vocabulary, not rewriting the
controller. The evaluation design enforces this division: the
generalisation protocol of Chapter~\ref{chap:eval} configures the
system without access to the held-out publication, which is only
meaningful because no publication-specific behaviour can hide inside
free-form controller output.
%% NOTE: this paragraph is the corrected version of the interim
%% sentence "The purpose is not to hard-code magazine-specific rules"
%% that the second marker flagged as self-contradictory; the fix is
%% the explicit tools-specific / controller-general division, with
%% transferability as the motivation, per the marker's own suggestion.

Within this contract, the controller's remaining discretion is real:
strategy selection per page, constraint composition per article, and
parameter adjustment within ranges are all consequential decisions,
and Chapter~\ref{chap:method} describes how they are made and logged.
What the contract removes is precisely the discretion that could not
be audited afterwards.
% =====================================================================
% Section 2.6 — Human-in-the-Loop and Constrained Translation
% (draft v1)
% Conventions as before. \ref placeholders: \ref{chap:method},
% \ref{sec:method-hitl}, \ref{chap:eval}, \ref{sec:background-agentic}.
% New bib keys introduced here (add to references.bib):
%   hokamp2017gbs   — Hokamp & Liu, ACL 2017, grid beam search.
%   dinu2019terminology — Dinu et al., ACL 2019, training-time
%                     terminology via inline annotation. %% VERIFY key
%                     details before adding; drop if 2.6 is over length.
% Existing keys reused: hasler2018terminology, yang2023human.
% This section grounds the interim claim "modeling human requirements
% as explicit constraints" that the second marker flagged as appearing
% from nowhere.
% =====================================================================

\section{Human-in-the-Loop and Constrained Translation}
\label{sec:background-hitl}

Magazine content is dense with items whose translation is a matter of
editorial decision rather than linguistic inference: brand and product
names, personal and place names with established renderings, recurring
column titles, and terms the publication deliberately leaves in the
source language. A translation workflow for publication therefore
needs a way for humans to constrain the machine --- and, just as
importantly, a defined account of what happens when human and machine
disagree. This section reviews how existing work represents and
enforces such constraints, and identifies the part of the problem that
remains open.

\subsection{Constrained Machine Translation}

Imposing pre-specified lexical constraints on neural translation has
been studied along two lines. The first enforces constraints
\emph{at decoding time}. Grid beam search extends beam search with an
additional dimension tracking how many constraints each hypothesis has
satisfied, guaranteeing that supplied words or phrases appear in the
output \cite{hokamp2017gbs}; Hasler et al.\ develop this direction
with finite-state tracking of constraint coverage, and frame adherence
to user-provided terminology as a recognised open challenge for neural
translation \cite{hasler2018terminology}, reporting gains of up to
3 BLEU from decoding with up to two dictionary constraints per
sentence.
%% VERIFIED V15 2026-08-13 https://aclanthology.org/N18-2081/ (Sec. 2.1
%% finite-state acceptors + multi-stack decoding, Sec. 2.2 attention-informed
%% constraint placement; Sec. 4.2 "Decoding with up to two dictionary
%% constraints per sentence yields gains of up to 3 BLEU". GBS
%% characterisation cross-checked against Hokamp & Liu 2017, Fig. 3, where
%% the added grid dimension is one row per constraint token.)
Constrained decoding offers hard guarantees, but the literature
records its costs. Decoding slows substantially: grid beam search
multiplies the number of beams by the number of constraints, giving a
runtime linear in that number rather than constant
\cite{hokamp2017gbs}, and Hasler et al.\ measure constrained decoding
running at between a fifth and a ninth of unconstrained speed for two
to four constraints per sentence \cite{hasler2018terminology}. And
forcing constraint tokens into the beam tends to reduce the fluency of
the surrounding translation: Post and Vilar show a decoder pushed to
place a lexically ill-fitting constraint distorting the sentence
around it, and attribute this to the model, once forced off the path
it prefers, no longer generating good competing hypotheses
\cite{post2018dba}.
%% VERIFIED V16 2026-08-14 speed half: https://aclanthology.org/P17-1141/
%% (Sec. 3.3, O(ktc) vs O(kt)) and https://aclanthology.org/N18-2081/
%% (Tab. 3, speed ratios 0.20/0.14/0.11 for c=2/3/4). Fluency half:
%% https://aclanthology.org/N18-1119/ (Fig. 6 caption and Sec. 6.4
%% "reference aversion"). The Post & Vilar evidence is a qualitative
%% demonstration plus model-score analysis, not a fluency score; the
%% hu2019vdba fallback was therefore not needed --- see the log.
Constraints in this line may be \emph{positive} (a rendering that must
appear) or \emph{negative} (a rendering that must not) --- a
distinction that maps directly onto the terminology choices and
do-not-translate rules of the present task.

The second line enforces constraints \emph{through the input}:
training-time or prompt-time methods surface the desired terminology
to the model --- for instance by annotating source terms inline with
their required translations \cite{dinu2019terminology} --- and rely on
the model to comply. This is the natural mode for LLM-based
translation, where a glossary can simply be included in the prompt,
and it preserves speed and fluency; but it trades away the guarantee.
Compliance becomes probabilistic, and nothing in the mechanism itself
reports whether a given constraint was honoured.

\subsection{Human-in-the-Loop Translation with LLMs}

LLMs widen the channel through which humans can steer translation
beyond term lists: Yang et al.\ show that revision instructions and
human feedback can guide an LLM translator towards more customised
output across successive passes; running GPT-3.5-turbo over five
domain-specific German--English benchmarks, they report a
domain-dependent improvement over direct prompting, for example 35.2
against 34.4 BLEU in the IT domain, with mixed results: in their
experiments only the comparison-based feedback variant consistently
improved over the baseline, while $k$-shot instruction variants often
scored below it \cite{yang2023human}. That free-form feedback helps
only under particular structures is itself instructive: it suggests
the bottleneck is not the channel but the governance of what flows
through it.
%% VERIFIED V17 2026-08-14 https://aclanthology.org/2023.mtsummit-users.8/
%% (Sec. 3 setup; Tab. 1 --- mixed-results wording per owner adjudication)
This restores the interactive character that end-to-end neural systems
had removed, and it matches the editorial reality of magazine work,
where preferences emerge from reviewing drafts rather than being fully
known in advance. But it also sharpens the enforcement problem of the
input-side line: instructions accumulated in prompts are free text,
their interaction is implicit, and whether any one of them was
followed on any given block is not directly observable.

\subsection{The Open Part: Adjudication, Priority, and Verifiability}

Read together, the two lines above address \emph{how} a known
constraint reaches the model, and the HITL work addresses \emph{how
humans express} preferences. What remains largely unaddressed is the
layer this project needs most: where the constraint set itself comes
from and how it is governed. Automatic terminology extraction ---
including the glossary construction built into the base system
(Section~\ref{sec:background-babeldoc}) --- produces candidate term pairs at
scale, but of varying quality; human editors hold authority but cannot
review every candidate on every run. Three questions follow that none
of the reviewed work answers: when a machine-extracted entry and a
human decision conflict, which wins, and by what rule; how a human
decision, once made, persists so that it is never re-litigated by the
machine on a later run; and how adherence is \emph{verified} on the
rendered output rather than assumed from the enforcement mechanism.

The position this project takes, developed in
Chapter~\ref{chap:method}, is that human decisions are sovereign
ground truth: machines draft candidates and machines validate
formats, but humans decide, through a two-pass review protocol whose
outcome is recorded as explicit, versioned data rather than prompt
text. Machine-extracted glossary entries are subordinated to these
decisions by construction --- conflicting automatic entries are
removed at glossary build time, not merely outranked at prompt time
--- so the priority rule is enforced at the point where the constraint
set is assembled rather than left to the model's discretion.
%% NOTE: "enforced where the set is assembled, not left to model
%% discretion" is the weakest-enforcement-point principle applied to
%% terminology; keep this phrasing aligned with the method chapter's
%% description of the glossary rebuild mechanism.
Because every constraint is explicit data with scope and priority, it
remains checkable after rendering: instruction adherence is an
evaluated property of the output (Chapter~\ref{chap:eval}), not a
presumed property of the mechanism. Representing human requirements as
explicit constraints is thus not an invention of this project --- it
is the founding idea of the constrained-MT literature above --- but
extending that idea from decode-time token constraints to
workflow-level, human-adjudicated, post-hoc verifiable decisions is
where this project's contribution in this area lies.
% =====================================================================
% Section 2.7 — Evaluating Translated Documents (draft v1)
% Same conventions as sections 2.3/2.4:
%  - %% VERIFY marks facts/formulas to re-check against primary sources.
%  - \ref placeholders: \ref{chap:eval}, \ref{sec:eval-midbreak},
%    \ref{sec:eval-metrics}, \ref{sec:corpus}, \ref{sec:background-hitl}.
%  - New bib keys introduced here (add to references.bib):
%    jiang2022blonde, lyu2021ltcr, wong2012cohesion, tan2022dcoem,
%    mueller2018contrapro, kikuchi2021layoutganpp, li2020layoutgan,
%    freitag2021mqm. Existing keys reused: papineni2002bleu,
%    rei2020comet, burchardt2013mqm, kocmi2023gemba_mqm,
%    lan2024translatotron.
% =====================================================================

\section{Evaluating Translated Documents}
\label{sec:background-eval}

Evaluation is central to this project because the usefulness of a
translated magazine depends jointly on two properties that fail
independently: whether the translation itself is adequate and coherent,
and whether the rendered page remains readable and visually faithful.
Evaluation is also unusually \emph{difficult} here, because the output
is a rendered PDF rather than plain text: a fluent translation can
still fail as a document by omitting blocks, leaving source text
behind, breaking article flow, or overlapping images, and none of these
failures is visible to text-only metrics. This section reviews how
existing work measures each property --- text quality at sentence
level, coherence at document level, error analysis, and layout quality
--- and identifies the combination this project adopts.

\subsection{Sentence-Level Text Metrics}

The standard automatic metrics for machine translation are
reference-based. BLEU \cite{papineni2002bleu} scores a candidate
against one or more reference translations by modified $n$-gram
precision with a brevity penalty:
\begin{equation}
\mathrm{BLEU} = \mathrm{BP} \cdot \exp\!\Big(\sum_{n=1}^{N} w_n \log p_n\Big),
\qquad
\mathrm{BP} =
\begin{cases}
1 & c > r\\
e^{\,1-r/c} & c \le r,
\end{cases}
\label{eq:bleu}
\end{equation}
where $p_n$ is the modified $n$-gram precision, $w_n$ the (typically
uniform) weights, and $c$, $r$ the candidate and reference lengths.
COMET \cite{rei2020comet} instead trains a neural regressor over
source, candidate, and reference embeddings to predict human quality
judgements, and correlates substantially better with human assessment
than surface-overlap metrics: on the WMT19 Metrics DARR corpus its
ranking model reaches a Kendall's $\tau$-like correlation of 0.449 on
English--Chinese against 0.235 for BLEU, and 0.445 on Chinese--English
against 0.321.
%% VERIFIED V18 2026-08-13 https://aclanthology.org/2020.emnlp-main.213/
%% (Tables 1 and 2, COMET-RANK rows)

Both metrics presuppose exactly what magazine translation cannot
guarantee: a clean, sentence-aligned text pair. Neither observes
layout, rendering, or coverage, and both are known to be insensitive to
document-level phenomena, which motivates the next family.

\subsection{Document-Level and Discourse-Aware Metrics}
\label{sec:background-doclevel}

Whether a translated article still reads as a continuous, coherent
text --- across paragraphs, columns, and pages --- is a document-level
property. Sentence-level metrics cannot distinguish document-level
improvements from sentence-level ones, a limitation identified
explicitly by Jiang et al.\ in proposing BlonDe
\cite{jiang2022blonde}: BlonDe categorises discourse-related spans
(entities, pronouns, tense, and discourse markers), computes a
similarity-based F1 over the categories, and in a large-scale human
study achieves significantly higher Pearson correlation with human
judgements than sentence-level metrics on BWB, a document-level
Chinese--English corpus \cite{jiang2022blonde}. BlonDe is directly
relevant to this project because its motivating language pair is
Chinese--English; however, its span categories were designed for
English-side evaluation, where tense morphology and explicit pronouns
carry the discourse signal. In the English-to-Chinese direction the
target language drops pronouns and does not inflect tense, so applying
BlonDe unmodified would measure the wrong surface phenomena; this
project therefore uses it only as a diagnostic, not a headline metric.
%% VERIFIED V19 2026-08-13 https://aclanthology.org/2022.naacl-main.111/
%% (Sec. 2: BWB is Chinese web novels with their English translations, so the
%% evaluated direction is zh->en; Sec. 5.3/Tab. 5: BlonDe obtains the highest
%% Pearson correlation at both units and significantly outperforms all other
%% metrics at document level, e.g. BLONDE.F1 .417 adequacy vs BLEU .343)

A simpler document-level extension --- d-BLEU, which matches
$n$-grams over the whole document rather than sentence by sentence
\cite{liu2020mbart} --- provides a cheap baseline but inherits the
reference-alignment assumption.
%% VERIFIED V20 2026-08-14 https://aclanthology.org/2020.tacl-1.47/ (n. 2:
%% "Standard BLEU scores match n-grams at sentence-level. We also consider
%% document-level where we match n-grams over the whole document").
%% The character-level document variant was deleted per owner decision:
%% no identifiable originating work.
A separate strand measures \emph{consistency} rather than similarity:
lexical cohesion was shown by Wong and Kit to distinguish human from
machine translation and, when added to sentence-level metrics, to
improve their correlation with human judgement by roughly 3--5
percentage points \cite{wong2012cohesion}; the lexical translation
consistency ratio (LTCR) \cite{lyu2021ltcr} formalises the intuition
directly, measuring whether repeated source terms are translated
consistently across a document. For a source word $w$ occurring $k$
times in a document with word-aligned translations $(t_1,\dots,t_k)$,
\begin{equation}
\mathrm{LTCR}(w) =
\frac{\sum_{i=1}^{k}\sum_{j=i+1}^{k} \mathbf{1}\!\left(t_i = t_j\right)}
     {C_k^2} \times 100\%,
\label{eq:ltcr}
\end{equation}
where $C_k^2$ is the size of the combination of the translation set and
$\mathbf{1}(t_i = t_j)$ returns 1 when the two translations are
identical; the metric extends to a corpus by summing numerators and
denominators over all words of interest \cite{lyu2021ltcr}.
%% VERIFIED V21 2026-08-13 https://aclanthology.org/2021.emnlp-main.262/
%% (Eq. 1, transcribed verbatim; LTCR is defined in that paper. NB the
%% venue is EMNLP 2021 main, not Findings --- bib corrected.)
LTCR matters for this project because article-level context injection
was introduced precisely to stabilise proper-name and terminology
rendering across pages (Chapter~\ref{chap:method}); LTCR turns that
design goal into a reportable number.

A third strand evaluates targeted phenomena through purpose-built test
suites rather than corpus-level scores: ContraPro tests context-aware
pronoun translation through contrastive example pairs
\cite{mueller2018contrapro}, and DCoEM contributes a test suite
covering four cohesive devices --- reference, conjunction,
substitution, and lexical cohesion \cite{tan2022dcoem}. This project
adopts the test-suite methodology rather than any specific suite: the
evaluation corpus (Section~\ref{sec:corpus}) contains real mid-sentence
page breaks identified and adjudicated during corpus preparation, and
each such break point functions as a targeted test case for exactly the
failure mode this project addresses. The resulting protocol ---
contrastive evaluation at splice points between page-scoped and
article-scoped translation --- is defined formally in
Chapter~\ref{chap:eval}.

\subsection{Error Taxonomies and Structured Human Evaluation}

Where automatic metrics compress quality into a scalar, the
Multidimensional Quality Metrics framework (MQM) treats translation
quality as a set of typed issues \cite{burchardt2013mqm}. MQM,
developed in the QTLaunchPad project by Lommel, Burchardt and
Uszkoreit, defines a hierarchy of issue types --- accuracy errors such
as omission and mistranslation, fluency errors, terminology errors,
and style issues --- from which task-appropriate subsets are selected,
with each identified error assigned a severity; quality is then an
aggregation over annotated errors rather than a single holistic score
\cite{burchardt2013mqm}. MQM-style annotation has since become the
reference standard for careful human MT evaluation: in the largest
such study to date, professional annotators working with full document
context produced a substantially different ranking of the WMT 2020
systems from the one established by crowd workers, with a clear
preference for human over machine output \cite{freitag2021mqm}.
%% VERIFIED V22 2026-08-13 https://aclanthology.org/2021.tacl-1.87/ (abstract)

GEMBA-MQM \cite{kocmi2023gemba_mqm} automates this style of
annotation: GPT-4 is prompted, with fixed few-shot examples, to mark
error spans and assign MQM categories and severities. In the WMT23
metrics evaluation it achieved state-of-the-art system-ranking
accuracy, at 96.5\% system-level pairwise accuracy in the preliminary
results its authors report. Its authors' stated caution is the
dependence on a proprietary, black-box model; this project mitigates
the reproducibility half of that concern by caching all judge
responses for exact replay and pinning the model version
(Chapter~\ref{chap:eval}).
%% VERIFIED V23 2026-08-14 https://aclanthology.org/2023.wmt-1.64/ (Tab. 1
%% for the accuracy; abstract, Conclusion and Limitations for the caution,
%% which concerns proprietary black-box dependence)
For this project the relevance is twofold: MQM supplies the issue
taxonomy for the structured human review rubric, and GEMBA-MQM
supplies a practical instrument for annotating errors \emph{at
specific locations} --- in particular at page-break splice points,
where the question is not ``how good is this document overall'' but
``did stitching these fragments introduce an omission, a
mistranslation, or a fluency break here''. Used this way, an LLM-based
annotator is a targeted microscope rather than a wholesale replacement
for human judgement, and its outputs are cached and replayable for
reproducibility (Chapter~\ref{chap:eval}).

\subsection{Layout and Rendering Quality}
\label{sec:background-layoutmetrics}

Text metrics of any granularity say nothing about whether the rendered
page is readable. Two bodies of work supply spatial measures. In
document image translation, Translatotron-V introduces
Structure-BLEU, which computes BLEU only over text pairs whose
regions correspond spatially (bounding-box IoU $\geq 0.5$), so that
translated text scores only where it also appears in the right
place --- coupling translation quality to spatial fidelity
\cite{lan2024translatotron}. This establishes
that layout preservation is a measurable property of translation
output rather than a purely aesthetic one.
%% VERIFIED V24 2026-08-14 https://arxiv.org/html/2407.02894 (metric
%% definition: OCR both images, match by bounding-box IoU, discard pairs
%% below 0.5, then compute BLEU over the remaining paired texts)

The layout generation literature provides the standard geometric
instruments; two are used throughout that field as measures of
perceptual layout quality, introduced by Li et al.\
\cite{li2020layoutgan} and normalised per element by Kikuchi et
al.\ \cite{kikuchi2021layoutganpp}:
\emph{Overlap}, the total pairwise overlapping area between element
bounding boxes on a page (normalised per element or per page), on the
principle that well-set elements do not intrude on one another; and
\emph{Alignment}, which for each element measures its minimum deviation
from alignment (left, centre, or right edges) with other elements, on
the principle that professionally set pages align elements to shared
axes.
\begin{equation}
\mathrm{Overlap} = \frac{1}{N}\sum_{i=1}^{N}\sum_{j \ne i}
\frac{A\!\left(b_i \cap b_j\right)}{A\!\left(b_i\right)}
\label{eq:overlap}
\end{equation}
where $b_i$ is the bounding box of element $i$ and $A(\cdot)$ its
area; the summation is Li et al.'s overlapping loss, and the leading
$1/N$ is Kikuchi et al.'s per-element normalisation of it. The
companion measure is defined over the same coordinates,
$(x^{L}, y^{T}, x^{C}, y^{C}, x^{R}, y^{B})$ being the top-left,
centre, and bottom-right coordinates of an element's box:
\begin{equation}
\mathrm{Alignment} = \frac{1}{N}\sum_{i=1}^{N}
\min\!\Big(
g\!\left(\Delta x_i^{L}\right),\,
g\!\left(\Delta x_i^{C}\right),\,
g\!\left(\Delta x_i^{R}\right),\,
g\!\left(\Delta y_i^{T}\right),\,
g\!\left(\Delta y_i^{C}\right),\,
g\!\left(\Delta y_i^{B}\right)
\Big),
\label{eq:alignment}
\end{equation}
where $g(x) = -\log(1-x)$ and
$\Delta x_i^{*} = \min_{\forall j \ne i}\left| x_i^{*} - x_j^{*}\right|$
for $* \in \{L, C, R\}$, with $\Delta y_i^{*}$ ($* \in \{T, C, B\}$)
computed in the same way. Equations~(10) and~(11) of Li et al.\
\cite{li2020layoutgan} state this sum without the leading $1/N$; the
per-element normalisation is Kikuchi et al.'s modification
\cite{kikuchi2021layoutganpp}.
%% VERIFIED V25 2026-08-14 https://arxiv.org/pdf/2009.05284 (Eqs. 9--11,
%% transcribed) + https://doi.org/10.1145/3474085.3475497 (Sec. 4.2.3: "We
%% use the Alignment and Overlap metrics used in the previous work [18].
%% We modify the original metrics by normalizing with the number of
%% elements N.")
These measures are appropriate for magazine translation for a reason
worth stating: they are \emph{reference-free} with respect to text.
They compare geometry against geometric principles (or against the
source page's own geometry), so they remain valid even where no
reference translation exists --- which is the normal case for magazine
content. BabelDoc's own evaluation already uses bounding-box IoU
against the source layout \cite{qi2026babeldoc}
(Section~\ref{sec:background-babeldoc}), and this project extends the same
geometric philosophy with task-specific measures --- among them a
mid-unit page-break rate and a block conservation invariant, defined
formally in Chapter~\ref{chap:eval} --- whose purpose is to detect the
failure modes that generic geometry misses.

\subsection{References, Their Limits, and Generalisation}

Two methodological choices complete the evaluation design. First, the
availability of an official human-translated Chinese edition of the
primary corpus magazine provides a rare parallel reference
(Section~\ref{sec:corpus}); but it is an \emph{editorial adaptation}
--- retranslated, re-set, and sometimes restructured by human editors
--- not a literal rendering of the English edition. Reference-based
scores computed against it are therefore treated as diagnostics, and
its principal evaluative use is comparative: where the official
edition itself resorts to a layout compromise (such as a mid-unit
break) that the system under test avoids, that comparison is evidence
about the difficulty of the problem, not merely about the system
(Section~\ref{sec:eval-midbreak}). Second, because the project claims
generalised magazine translation rather than adaptation to a single
publication, evaluation uses leave-one-publication-out cross-validation:
all adaptive components are configured without access to the held-out
publication, and performance on it measures transfer rather than fit.
The corresponding protocol, and the anti-overfitting constraint that
forbids publication-specific rules anywhere in the pipeline, are
specified in Chapter~\ref{chap:eval}.

In combination, the evaluation instruments adopted from this
literature form three layers, each targeting a property the others
miss: geometric and structural measures for the rendered page
(Section~\ref{sec:background-layoutmetrics}), contrastive splice-point
evaluation with MQM-style annotation for translation quality at the
exact locations this project's mechanisms change
(Section~\ref{sec:background-doclevel}), and document-level consistency
measures for article-scoped coherence. Their formal definitions, and
the conservation invariants that guard the pipeline against silent
loss, are given in Chapter~\ref{chap:eval}.
